\documentclass[11pt]{article}
\usepackage{preamble}
\renewcommand{\answer}[1]{}

\begin{document}

\ladderheading{AP Statistics \textperiodcentered\ Unit 3 \textperiodcentered\ Topic 3.12 \textperiodcentered\ B: 2026-12-15 \textperiodcentered\ E: 2027-01-27}{Two Analysts, One Planned Test}

\focusbanner{TODAY'S PROBLEM}

Suppose an online retailer has 50,000 active accounts. Before collecting data, its team asks, ``Does a progress meter change the proportion of active accounts completing checkout?'' It randomly selects 360 accounts and randomly assigns 180 to each screen. The table shows the outcomes. Let $p_M$ and $p_S$ be the corresponding population completion proportions.\par\smallskip \textbf{Nia:} ``Use $H_0:p_M-p_S=0$ versus $H_a:p_M-p_S\ne0$ because either direction is a change.''\par \textbf{Omar:} ``Because $\hat p_M$ is larger, use $H_a:p_M-p_S>0$.''\par
\begin{center}
\textbf{Checkout outcomes}\par\smallskip
\begin{tabular}{|l|c|c|c|}
\hline
\textbf{Screen} & \textbf{n} & \textbf{Complete} & \textbf{Other} \\ \hline
\textbf{Meter} & 180 & 126 & 54 \\ \hline
\textbf{Standard} & 180 & 108 & 72 \\ \hline
\end{tabular}
\end{center}
\begin{firsttakebox}{FIRST TAKE --- before the video (5 min)}
When the team asks whether checkout changes, should a decrease count too? Use the original question to explain.
\workspace{0.9in}
\end{firsttakebox}

\begin{callout}{\IconBook\ SET THE QUESTION, PARAMETERS, AND CONDITIONS (keep on the page)}{calloutgreen}
Define population proportions before writing hypotheses. For no difference,
$H_0:p_1-p_2=0$. Choose $H_a:p_1-p_2>0$, $<0$, or $\ne0$ from the research
question set \textbf{before} viewing results; do not choose a tail because of the observed sign.
Use a two-sample $z$-test for a difference of proportions with independent random samples
or randomized groups. Check each 10 percent condition when sampling without replacement.
Under $H_0$, pool: $\hat p_c=(x_1+x_2)/(n_1+n_2)$; all four pooled expected counts must be at least 10.
\end{callout}

\newpage
\textbf{Finish after the video:}\par\smallskip

\dokbadge{1}~\textbf{(a)} Define $p_M$ and $p_S$ in context. Calculate $\hat p_M$, $\hat p_S$, and the observed difference $\hat p_M-\hat p_S$.\par
\workspace{0.7in}

\dokbadge{2}~\textbf{(b)} Name the test procedure. Calculate the pooled proportion and verify randomization, the 10 percent condition, and all four pooled expected counts.\par
\workspace{1.4in}

\dokbadge{3}~\textbf{(c)} In 3--4 sentences, choose and state the hypotheses. Use the original word change and when the direction was chosen. Explain how Omar's choice after seeing the result changes the question and could overstate the evidence.\par
\workspace{2.0in}

\begin{sentenceframebox}\raggedright\textbf{Frame:} The alternative is \blank[1.6in] because the original question asks \blank[2.0in].\end{sentenceframebox}

\begin{sentenceframebox}\raggedright\textbf{Frame:} Choosing the direction after seeing \blank[1.2in] could mislead readers by \blank[2.5in].\end{sentenceframebox}

\begin{summaryexitbox}{EXIT --- turn this in}
One thing the video changed about my first take: \hrulefill\par
\workspace{0.4in}
\end{summaryexitbox}

\end{document}
